\documentclass[11pt]{article}

% ---------------------------------------------------------------------------
% C3i Hub IoT Security Testbed — technical report (LaTeX source)
% Compile: pdflatex C3i_IoT_Testbed_Report.tex  (run twice for the ToC)
% ---------------------------------------------------------------------------
\usepackage[T1]{fontenc}
\usepackage[utf8]{inputenc}
\usepackage[a4paper,margin=0.95in]{geometry}
\usepackage[protrusion=true,expansion=false]{microtype}
\usepackage{parskip}
\usepackage{xcolor}
\usepackage{graphicx}
\usepackage{booktabs}
\usepackage{tabularx}
\usepackage{array}
\usepackage{enumitem}
\usepackage[most]{tcolorbox}
\usepackage{titlesec}
\usepackage{fancyhdr}
\usepackage[hidelinks]{hyperref}

% ---- brand palette --------------------------------------------------------
\definecolor{accent}{HTML}{149C77}   % mint-teal, darkened for print contrast
\definecolor{accentbg}{HTML}{E7F6F0}
\definecolor{ink}{HTML}{11181E}
\definecolor{slate}{HTML}{3F4B57}
\definecolor{rulegray}{HTML}{D4DBE1}
\definecolor{warnbar}{HTML}{C07A1A}
\definecolor{warnbg}{HTML}{FBF2E2}
\definecolor{codebg}{HTML}{F2F5F7}

\hypersetup{colorlinks=true, linkcolor=accent, urlcolor=accent, citecolor=accent}

% ---- headings -------------------------------------------------------------
\titleformat{\section}{\Large\bfseries\color{ink}}{\thesection\quad}{0pt}{}[{\color{accent}\titlerule[1.2pt]}]
\titleformat{\subsection}{\large\bfseries\color{ink}}{\thesubsection\quad}{0pt}{}
\titlespacing*{\section}{0pt}{16pt}{8pt}
\titlespacing*{\subsection}{0pt}{11pt}{4pt}

% ---- header / footer ------------------------------------------------------
\pagestyle{fancy}
\fancyhf{}
\renewcommand{\headrulewidth}{0pt}
\renewcommand{\footrulewidth}{0.4pt}
\fancyfoot[L]{\footnotesize\color{slate}C3i Hub · IoT Security Testbed}
\fancyfoot[R]{\footnotesize\color{slate}\thepage}

% ---- helpers --------------------------------------------------------------
% Inline code: boxed in prose, plain in tables (set \intable=1 there) to avoid
% non-breaking colorbox cells overflowing narrow columns.
\newif\ifintable\intablefalse
\newcommand{\code}[1]{\ifintable\texttt{#1}\else{\small\colorbox{codebg}{\texttt{#1}}}\fi}
\newcolumntype{L}[1]{>{\raggedright\arraybackslash}p{#1}}

\newtcolorbox{callout}[1][]{colback=accentbg,colframe=accent,boxrule=0pt,
  leftrule=3pt,arc=1pt,left=8pt,right=8pt,top=6pt,bottom=6pt,#1}
\newtcolorbox{warnbox}[1][]{colback=warnbg,colframe=warnbar,boxrule=0pt,
  leftrule=3pt,arc=1pt,left=8pt,right=8pt,top=6pt,bottom=6pt,#1}
% Full-width wrapping code display (breaks at spaces; no overflow).
\newtcolorbox{codebox}{colback=codebg,boxrule=0pt,arc=1pt,
  left=10pt,right=10pt,top=6pt,bottom=6pt,before skip=7pt,after skip=7pt}

\setlist[itemize]{leftmargin=1.2em,itemsep=2pt,topsep=3pt}
\setlist[enumerate]{leftmargin=1.5em,itemsep=2pt,topsep=3pt}
\renewcommand{\arraystretch}{1.3}

% ===========================================================================
\begin{document}
\sloppy

% ---- cover ----------------------------------------------------------------
\begin{titlepage}
\thispagestyle{empty}
\vspace*{3cm}
\noindent{\color{accent}\rule{3cm}{3pt}}\\[1.2cm]
{\Huge\bfseries\color{ink} C3i Hub IoT Security Testbed}\\[0.5cm]
{\LARGE\color{slate} What the platform does, how to run it,\\ and how to verify its results}\\[1.5cm]
\noindent{\large A unified Bluetooth/BLE, Wi-Fi/IP, hardware (UART/SPI/I2C) and
firmware auditing platform developed under the \emph{Development of Testbeds and
Security Mechanisms for Smart IoT Devices} initiative.}\\[2cm]
\noindent\begin{tabular}{@{}ll@{}}
{\bfseries Version} & Platform v2.0.0\\
{\bfseries Domains} & Bluetooth/BLE, Wi-Fi/IP, Hardware, Firmware\\
{\bfseries Scope} & Authorized, evidence-backed IoT security assessment\\
{\bfseries Source} & \texttt{docs/C3i\_IoT\_Testbed\_Report.tex} (this document)\\
\end{tabular}
\vfill
\noindent{\color{accent}\rule{\textwidth}{0.6pt}}\\[2pt]
{\footnotesize\color{slate} C3i Hub · IIT Kanpur \hfill Technical report}
\end{titlepage}

% ---- contents -------------------------------------------------------------
\tableofcontents
\newpage

% ===========================================================================
\section{Executive summary — is this tool trustworthy?}
The C3i IoT Security Testbed is a single workspace that assesses a connected
device across four layers: \textbf{Bluetooth/BLE}, \textbf{Wi-Fi/IP}, a
\textbf{hardware testbed} (UART, SPI, I2C, logic analysis and bounded fault
injection) and \textbf{firmware static analysis}. It is built around three
principles that make its output defensible in a client report:

\begin{itemize}
  \item \textbf{Honesty over theatre.} Every module probes for the capabilities
  it needs and reports \emph{unavailable} rather than fabricating a result. A
  distance is a \emph{model estimate}, a CVE is a \emph{candidate}, and a hash
  identifies the bytes saved locally --- never a claim the tool cannot support.
  \item \textbf{Authorization first.} Non-passive actions are gated behind a
  cryptographically signed Rules of Engagement (RoE) \emph{and} a host-level
  invasive switch. Scope is enforced per target.
  \item \textbf{Evidence at every step.} Each run records logs, artefacts
  (hashed), and structured findings that roll up into an engagement report.
\end{itemize}

\begin{callout}
\textbf{Why ``is it sus?'' is the right question.} The platform is designed so
that a sceptical reviewer can trace every finding to its basis: the module, the
command issued, the raw evidence, and the confidence. Nothing is asserted that
the collected evidence does not support.
\end{callout}

% ===========================================================================
\section{Platform at a glance}
The workflow is linear and auditable:
\begin{center}
\begin{codebox}\small\ttfamily (1) Engagements $\rightarrow$ (2) Firmware / Hardware / Module runner $\rightarrow$ (3) Findings $\rightarrow$ (4) Evidence \& reports\end{codebox}
\end{center}

A FastAPI backend hosts a plugin registry (36 modules across four domains); a
React/TypeScript console drives it. Each module declares its parameters,
safety level, required capabilities and tags, so the UI and the authorization
gate are generated from the same source of truth.

\begin{table}[h]
\centering\footnotesize\intabletrue
\begin{tabularx}{\textwidth}{@{}L{3.1cm} L{4.2cm} X@{}}
\toprule
\textbf{Domain} & \textbf{Input} & \textbf{What it produces}\\
\midrule
Bluetooth/BLE & A device in radio range & Discovery, OS/version inference, firmware ID, GATT map, pairing/privacy posture, coarse distance, CVEs\\
Wi-Fi/IP & An owned host or AP & Service/version scan, TLS posture, wireless survey, MQTT audit, bounded deauth\\
Hardware & A board / debug interface & UART console + boot log, SPI/I2C identity \& dump, logic decode, complete filesystem acquisition, bounded glitch\\
Firmware & An image or recovered partition & Carving, entropy, secrets, SBOM, secure-boot audit, CVE mapping\\
\bottomrule
\end{tabularx}
\end{table}

The domains map to the layers of a device (radio, network, firmware, memory,
physical interface) so a client report can show coverage per layer.

\subsection{Architecture \& plugin model}
Each domain is a \emph{plugin}; each capability is a \emph{module} declared with
a decorator carrying its name, title, safety level, Pydantic parameter model,
required capabilities and tags. The registry exposes this metadata over the API,
so three things are generated from one declaration: the console's input form, the
authorization requirements, and the capability pre-flight. A \emph{capability}
(e.g.\ \code{net.tshark}, \code{bt.dbus}, \code{host.block\_tools}) is probed at
runtime; when absent, the module says so and offers the exact install command,
instead of failing opaquely or faking a result. Jobs run through a serialized
runner per hardware/radio domain (so two jobs never fight over one adapter,
serial port or bus), emitting live progress, logs, structured findings and
hashed evidence --- all persisted and attributable to the engagement.

\subsection{Intrusiveness ladder}
\begin{table}[h]
\centering\footnotesize\intabletrue
\begin{tabularx}{\textwidth}{@{}L{2.4cm} X@{}}
\toprule
\textbf{Level} & \textbf{Meaning \& examples}\\
\midrule
Passive & Observation only; always allowed. Discovery, GATT read, port scan, entropy.\\
Active & Touches the target but non-destructive. Baud detect, filesystem dump, default-credential test.\\
Invasive & State-changing / stress; bounded and reversible; double-gated. Voltage glitch, deauth, link stress, MAC spoof.\\
\bottomrule
\end{tabularx}
\end{table}

% ===========================================================================
\section{The safety model: why every finding is defensible}
Authorization is enforced in \code{assert\_authorized} before any non-passive
module runs. The gate checks, in order:
\begin{enumerate}
  \item \textbf{Host switch} --- invasive/active modules require
  \code{C3I\_ENABLE\_INVASIVE=true} on the host.
  \item \textbf{Engagement status} --- must be \emph{authorized} or \emph{active}.
  \item \textbf{Signed RoE} --- an HMAC-SHA256 signature over the engagement id,
  signer, scope and expiry. Changing the scope invalidates the signature
  (\emph{tamper-evident}).
  \item \textbf{Expiry} --- an expired RoE blocks execution.
  \item \textbf{Scope} --- the target (MAC, IP, CIDR, device path, or the
  pseudo-target \code{local}) must be in the authorized set.
\end{enumerate}

Modules are classified \textbf{passive} (observation only; always allowed),
\textbf{active} (touches the target but non-destructive) or \textbf{invasive}
(state-changing / stress; always reversible and bounded). The console surfaces
the level on every module and requires a pre-flight checklist for invasive runs.

% ===========================================================================
\section{Bluetooth / BLE}
The Bluetooth domain profiles a device, identifies its OS/firmware, enumerates
the full GATT attribute table and classifies access, assesses pairing/privacy
posture, verifies encryption, maps CVEs, and gives a coarse distance with honest
uncertainty. It drives a real BlueZ stack and degrades honestly when a
dependency or adapter is absent.

\begin{table}[h]
\centering\footnotesize\intabletrue
\begin{tabularx}{\textwidth}{@{}L{3.3cm} L{1.9cm} X@{}}
\toprule
\textbf{Module} & \textbf{Safety} & \textbf{What it does}\\
\midrule
\code{scan} & passive & Discovers classic + LE devices, ranks by RSSI, flags strong unnamed devices, adds vendor/OS hints\\
\code{profile} & passive & SDP/class + manufacturer-data evidence to infer OS family/version with a confidence score; flags pairing hygiene\\
\code{gatt\_firmware} & passive & Reads the Device Information Service (DIS) over D-Bus (gatttool fallback) and maps firmware to CVEs\\
\code{gatt\_enum} & passive & Walks the \emph{entire} GATT table; classifies each characteristic OPEN vs PROTECTED; bounded reads\\
\code{le\_security} & passive & Address privacy (public vs random), pairing/bonding state, services-without-pairing\\
\code{link\_rtt} & passive & L2CAP echo RTT + packet delivery; coarse RSSI distance estimate\\
\code{capture\_audit} & passive & Encrypted-vs-plaintext verdict from a capture + version hints + CVEs\\
\code{cve\_lookup} & passive & Maps OS/firmware keywords to CVEs\\
\code{link\_stress} & invasive & Bounded L2CAP echo load ($\le$1000 pkts); measures degradation, always stops\\
\code{mac\_spoof} & invasive & Changes the local adapter identity then restores it\\
\bottomrule
\end{tabularx}
\end{table}

\subsection{Telling how far away a device is}
\code{link\_rtt} reports an \emph{Estimated distance} from the device's current
RSSI using a log-distance path-loss model:
\begin{codebox}\small\ttfamily d = 10\textasciicircum{}((RSSI\_at\_1m $-$ RSSI) / (10 $\times$ n))\end{codebox}
\begin{itemize}
  \item \textbf{RSSI\_at\_1m} defaults to $-59$\,dBm --- a typical BLE 1\,m value,
  deliberately \emph{not} the advertised TxPower.
  \item \textbf{n} (path-loss exponent) defaults to 2.5 indoor (2.0 free space,
  $\sim$3.0 obstructed); the operator can override it.
  \item The card shows a point estimate, a plausible range (varying $n$ by
  $\pm0.5$), the RSSI, the reference and a low/medium confidence --- never
  ``high'' from a single sample. RTT is reported separately for link quality
  only; it is \emph{not} converted to distance.
\end{itemize}
\begin{warnbox}
\textbf{Why it is deliberately fuzzy.} Body shadowing, antenna orientation and
multipath routinely cause 2--3$\times$ error on one RSSI sample. For a tighter
answer, measure RSSI at 1\,m in the same environment and pass it as a reference.
For true ranging use UWB / time-of-flight --- Bluetooth RSSI is not a tape measure.
\end{warnbox}

\subsection{Deep GATT enumeration \& pairing/privacy (the core BLE pentest)}
Reading only the DIS is a fingerprint, not an assessment. \code{gatt\_enum}
walks every service, characteristic and descriptor over the BlueZ D-Bus API and
classifies each characteristic's access from its \texttt{Flags}:
\begin{itemize}
  \item \textbf{OPEN} --- reachable with \emph{no} encryption/authentication
  (plain \texttt{read} / \texttt{write} / \texttt{write-without-response}). An OPEN
  \emph{writable} characteristic on an actuator, lock, config or OTA service is a
  classic high-severity finding and is reported as such.
  \item \textbf{PROTECTED} --- requires a bonded/encrypted link
  (\texttt{encrypt-*} / \texttt{secure-*} / \texttt{authenticated-signed-writes}).
  \item OPEN-readable values are read (bounded) and saved as evidence so you can
  check for secrets or PII exposed without pairing.
\end{itemize}
\code{le\_security} reports the observable posture: a \emph{fixed public address}
(trackable --- flagged) vs a rotating random/resolvable one; paired/bonded/trusted
state; and services exposed \emph{without pairing}. It never invents an SMP method
it did not witness.
\begin{callout}
\textbf{Grounded in real BLE methodology} --- service/characteristic discovery
then OPEN-vs-PROTECTED review, and the spec-level weaknesses behind KNOB (minimum
key length), BIAS (impersonation) and Just-Works/legacy pairing (no MITM
protection; sniffable keys before LE Secure Connections). Enumeration is not an
auth bypass --- it only records what the device exposes unauthenticated.
\end{callout}

\subsection{What more can be done}
Active SMP probing (detect Just-Works/downgrade, behind the invasive gate);
passive advertisement sniffing with an nRF/Ubertooth; authorized
characteristic write-testing; explicit SweynTooth/BleedingTooth/KNOB/BIAS
mapping; and OTA/DFU service detection.

% ===========================================================================
\section{Wi-Fi / IP}
This domain assesses the network, session and OS layers of an \textbf{owned}
target or access point --- honest service and posture assessment, not cracking a
third party's network.

\begin{table}[h]
\centering\footnotesize\intabletrue
\begin{tabularx}{\textwidth}{@{}L{3.2cm} L{1.8cm} X@{}}
\toprule
\textbf{Module} & \textbf{Safety} & \textbf{What it does}\\
\midrule
\code{discovery} & passive & nmap host sweep to inventory the device mesh\\
\code{portscan} & passive & Service/version scan; flags dangerous services; version$\rightarrow$CVE\\
\code{os\_fingerprint} & passive & TTL heuristic + optional nmap -O\\
\code{tls\_audit} & passive & Enumerates TLS versions/ciphers; flags downgrade + self-signed\\
\code{wireless\_survey} & passive & Nearby APs, encryption, WPS exposure (iw)\\
\code{mqtt\_audit} & passive & Detects anonymous-subscribe brokers, samples topics\\
\code{traffic\_capture} & passive & Captures N seconds, verifies encryption, mines version hints + CVEs\\
\code{default\_credentials} & active & Bounded default-credential test on an HTTP UI\\
\code{deauth\_test} & invasive & $\le$20 deauth frames; detects missing 802.11w; restores the interface\\
\bottomrule
\end{tabularx}
\end{table}

A no-RF Wi-Fi lab (\code{mac80211\_hwsim} + hostapd) lets the radio-layer modules
run deterministically on one Linux host with no transmitter (see \S\ref{sec:testing}).

% ===========================================================================
\section{Hardware testbed --- UART / SPI / I2C and complete filesystem extraction}
This is the Raspberry-Pi rig domain, derived from the C3i Hub IOTSMS-2023
testbed. For a camera or similar device the goal is to get from a bare board to a
\textbf{complete, analysable filesystem} while staying read-only first and
preserving evidence.

\begin{table}[h]
\centering\footnotesize\intabletrue
\begin{tabularx}{\textwidth}{@{}L{3.4cm} L{1.6cm} X@{}}
\toprule
\textbf{Module} & \textbf{Safety} & \textbf{What it does}\\
\midrule
\code{hardware\_info} & passive & Rig diagnostics: GPIO backend, serial, probes, storage\\
\code{detect\_baud} & active & Auto-discovers UART baud by scoring boot output\\
\code{capture\_bootlog} & active & Power-cycles, captures + analyses the boot log, maps CVEs\\
\code{uart\_console} & active & Detects console/login/bootloader; optional safe U-Boot breakout\\
\code{serial\_exec} & active & Runs one bounded command on an existing shell\\
\code{filesystem\_dump} & active & Recovers one \texttt{/proc/mtd} partition --- or \emph{all} (\texttt{all\_partitions})\\
\code{rootfs\_archive} & active & Tars the \emph{live} mounted filesystem over UART or netcat\\
\code{host\_block\_acquire} & active & Read-only image of a USB/SD device attached to the host\\
\code{adb\_acquire} & active & Tars a filesystem off an Android device via adb exec-out\\
\code{spi\_identify} & active & Read-only JEDEC identity + write-protection state\\
\code{spi\_dump} & invasive & Reads the flash via flashrom; hands off to firmware\\
\code{i2c\_scan} & active & Scans an I2C bus / reads an EEPROM candidate\\
\code{logic\_capture} & active & Passive logic-analyzer capture + UART/SPI/I2C decode\\
\code{power\_control} & active & Manual relay / optocoupler on/off/cycle\\
\code{voltage\_glitch} & invasive & Bounded glitch pulse train; always restores power\\
\bottomrule
\end{tabularx}
\end{table}

\subsection{Complete-acquisition method matrix}
``Extract the whole filesystem'' has several correct answers depending on how the
device is wired and whether it is running. Each method writes a SHA-256-identified
image to evidence and hands off to the firmware domain.

\begin{table}[h]
\centering\footnotesize\intabletrue
\begin{tabularx}{\textwidth}{@{}L{4.3cm} L{3.6cm} X@{}}
\toprule
\textbf{Situation} & \textbf{Module} & \textbf{Why}\\
\midrule
Live, UART shell, NOR flash & \code{filesystem\_dump} all & Block-exact dump of every partition + manifest\\
Live, UART shell, NAND/UBI/overlay & \code{rootfs\_archive} (serial) & A block \texttt{dd} of a NAND mtdblock is meaningless; tar the mounted FS\\
Live device with an IP & \code{rootfs\_archive} (netcat) & DUT pushes tar over TCP; $100\times$ faster than UART base64\\
Storage attached to the PC (USB/SD) & \code{host\_block\_acquire} & Read-only image at USB speed; refuses the host's own disks\\
Android-based device & \code{adb\_acquire} & Binary-safe tar over adb exec-out\\
Bare SPI-NOR chip & \code{spi\_identify} $\rightarrow$ \code{spi\_dump} & flashrom reads the chip directly (invasive, gated)\\
\bottomrule
\end{tabularx}
\end{table}

\begin{warnbox}
\textbf{What acquisition is --- and is not.} None of these bypass a locked
bootloader, secure boot or an unknown login; they acquire what an existing shell
or an exposed storage interface already gives. \code{host\_block\_acquire} is
strictly read-only, never mounts the target, and refuses the host's own system
disks. Chip-off NAND and JTAG remain out of scope (no software path). Base64 over
115200 baud carries $\sim$8.4\,KiB/s, so a 16\,MiB partition takes $\sim$33\,min;
the netcat transport is far faster when the device has an IP.
\end{warnbox}

% ===========================================================================
\section{Firmware static analysis}
Whether an image arrives from a UART/SPI dump, a host-block image or a vendor
download, the firmware domain runs one pipeline:
\begin{codebox}\small\ttfamily signatures $\rightarrow$ entropy $\rightarrow$ strings $\rightarrow$ secrets $\rightarrow$ SBOM $\rightarrow$ CVE map $\rightarrow$ secure-boot audit\end{codebox}
\begin{itemize}
  \item \code{analyse} --- end-to-end carve, entropy, secrets, components, defaults, CVEs.
  \item \code{extract} --- binwalk or built-in bounded gzip/CPIO/tar/ZIP carving $\rightarrow$ ZIP.
  \item \code{entropy} --- per-block Shannon entropy to spot encryption/compression.
  \item \code{secrets} --- keys, cloud credentials, password assignments, shadow hashes (masked).
  \item \code{sbom} --- component inventory (BusyBox, OpenSSL, dropbear\dots) + CVE mapping.
  \item \code{boot\_chain} --- static secure-boot evidence (signed FIT/AVB/dm-verity, fuses, debug locks).
\end{itemize}
\begin{callout}
\textbf{Honesty rule.} A component version string is a CVE \emph{candidate}, not
proof of exploitability; \code{boot\_chain} reports gaps as \emph{not evidenced}.
Signature detection is not unpacking --- UBI/UBIFS, encrypted updates and
proprietary containers need their own tooling.
\end{callout}

% ===========================================================================
\section{Vulnerability intelligence}
CVE mapping draws on three public sources with an offline SQLite FTS5 cache, so
the platform works air-gapped once the cache is warmed.
\begin{table}[h]
\centering\footnotesize\intabletrue
\begin{tabularx}{\textwidth}{@{}L{3.2cm} L{4.6cm} X@{}}
\toprule
\textbf{Source} & \textbf{Role} & \textbf{Notes}\\
\midrule
NVD 2.0 API & Canonical CVE records, CVSS, CPE & Rate-limited 5/30\,s (50 with a key)\\
OSV.dev & Ecosystem advisories / patch info & Open API\\
EPSS (first.org) & Exploit-prediction probability + percentile & Free\\
Built-in cache & SQLite FTS5 offline store + risk scoring & \code{C3I\_CVE\_ONLINE=false} to go offline\\
\bottomrule
\end{tabularx}
\end{table}
Per-finding severity is normalised from CVSS (critical $\ge$9.0, high $\ge$7.0,
medium $\ge$4.0). Engagement risk is a diminishing-returns roll-up capped at 100.

% ===========================================================================
\section{Reporting \& evidence}
Every job persists structured findings (title, category, severity, CVSS/EPSS,
CVE/CWE ids, remediation, evidence) and hashed artefacts. An engagement rolls
these into a downloadable report. CVE matches are labelled \emph{candidate}, and
findings carry the basis of each conclusion so a reviewer can audit them.

% ===========================================================================
\section{Operating system: Kali vs Attify OS vs Debian / Pi OS}
\begin{table}[h]
\centering\footnotesize\intabletrue
\begin{tabularx}{\textwidth}{@{}L{3cm} X L{3.4cm}@{}}
\toprule
\textbf{Option} & \textbf{Fit} & \textbf{Verdict}\\
\midrule
Raspberry Pi OS (64-bit) & Native GPIO/SPI/I2C/UART on the Pi rig; stable & \textbf{Best for the hardware rig}\\
Kali Linux & Rich wireless/BLE tooling preinstalled; good on a laptop & Best for the radio/network workstation\\
Attify OS & IoT-focused bundle; now dated and unmaintained & Optional; cherry-pick tools instead\\
Debian/Ubuntu & Clean base; install only what you need & Fine for the backend server\\
\bottomrule
\end{tabularx}
\end{table}
\textbf{Recommendation:} Pi OS on the bench controller for the hardware domain,
and Kali (or Debian + the listed tools) on the laptop that runs the radio and
network modules. The capability probe tells you exactly which tools are missing
and gives the \texttt{apt}/\texttt{pip} command to install each.

% ===========================================================================
\section{Do you need a Raspberry Pi?}
\textbf{No, not for everything.} Firmware analysis, Wi-Fi/IP, Bluetooth, and
USB-based UART/SPI all run on a laptop. A Raspberry Pi 5 is recommended when you
want \emph{bench automation}: GPIO-driven power control, relay/optocoupler
switching, glitch injection, and native SPI/I2C headers. Wire the device
\textbf{powered off}, power it through the relay/optocoupler (never the Pi
header), and keep a common ground. Full pin map and bill of materials are in
\texttt{docs/HARDWARE\_SETUP.md}.

% ===========================================================================
\section{Running the tool}
\begin{enumerate}
  \item Start the backend (add \code{C3I\_ENABLE\_INVASIVE=true} to allow
  non-passive modules), then the console, then open \url{http://localhost:5173}:
  \begin{codebox}\small\ttfamily
  cd backend \&\& ../.venv/bin/uvicorn app.main:app --port 8080\\
  cd web \&\& npm run dev
  \end{codebox}
  \item Sign in, create an engagement, sign its RoE (set the scope), then run
  modules from the runner or the domain pages.
\end{enumerate}
On a laptop the hardware rig operations report the mock backend and say so; the
software still runs end to end.

% ===========================================================================
\section{Testing \& labs}
\label{sec:testing}
Everything is testable offline:
\begin{itemize}
  \item \textbf{IP-service simulator} --- a loopback target exercises
  \code{discovery}, \code{portscan}, \code{tls\_audit}, \code{default\_credentials},
  \code{mqtt\_audit} with no radio.
  \item \textbf{No-RF Wi-Fi lab} --- \code{mac80211\_hwsim} + hostapd stand up a
  virtual AP + client so \code{wireless\_survey}, capture and \code{deauth\_test}
  run deterministically (\texttt{docs/WIFI\_HWSIM\_LAB.md}).
  \item \textbf{Emulated camera console} --- the test suite's \texttt{FakeCamDevice}
  speaks busybox over a fake serial line to exercise the UART$\rightarrow$filesystem
  path (login, \texttt{/proc/mtd}, \texttt{dd|base64}, \texttt{tar}).
  \item \textbf{Acquisition tests} --- multi-partition dump-all, live-rootfs tar,
  a real loopback socket for netcat exfil, a monkeypatched \texttt{dd} for
  host-block imaging (incl. system-disk refusal), and a stubbed adb.
\end{itemize}
Run the suite with \code{cd backend \&\& ../.venv/bin/pytest -q}.

% ===========================================================================
\section{Simulations: router, Bluetooth and Wi-Fi without hardware}
Because real RF and real boards are not always available, the platform supports
deterministic, no-hardware stand-ins so workflows can be built, demonstrated and
regression-tested safely.
\begin{itemize}
  \item \textbf{Router / IP device} --- the loopback \emph{IP-service simulator}
  presents HTTP, TLS and MQTT endpoints so \code{discovery}, \code{portscan},
  \code{tls\_audit}, \code{default\_credentials} and \code{mqtt\_audit} produce
  real findings against a controlled target.
  \item \textbf{Wi-Fi} --- \code{mac80211\_hwsim} creates virtual 802.11 radios
  driven by hostapd and wpa\_supplicant. A virtual AP (deliberately weak posture:
  PMF off) and an associated client let \code{wireless\_survey}, monitor-mode
  capture and \code{deauth\_test} run against a station that actually
  de-authenticates and re-associates --- with no transmitter and no spectrum use.
  \item \textbf{Bluetooth} --- the \code{FakeCamDevice} serial console and
  canned BlueZ/D-Bus maps exercise the console-to-filesystem path and the GATT
  classifier offline; a real BlueZ adapter is used when present.
\end{itemize}
These simulations are the same targets the automated test suite uses, so a
demonstration and a regression run share one source of truth.

% ===========================================================================
\section{Verifying results against real test vectors}
\begin{itemize}
  \item \textbf{Firmware carving} --- run \code{analyse}/\code{extract} on a known
  OpenWrt/vendor image and confirm the carved SquashFS/partition offsets match
  \texttt{binwalk} run independently.
  \item \textbf{Acquisition integrity} --- compare the module's SHA-256 against an
  independent \texttt{sha256sum} of the same source.
  \item \textbf{CVE mapping} --- confirm a component/version maps to the expected
  CVE ids in NVD, and that EPSS/CVSS values match the public record.
  \item \textbf{GATT access} --- cross-check OPEN/PROTECTED classification against
  the device's advertised characteristic flags in nRF Connect.
\end{itemize}

% ===========================================================================
\section{Limitations}
\begin{itemize}
  \item Does not bypass locked bootloaders, secure boot or unknown credentials.
  \item RSSI distance is a model estimate, not a measurement.
  \item CVE matches are candidates from version strings, not proof of exploitability.
  \item Chip-off NAND, JTAG/SWD and RF jamming are out of scope by design.
  \item Radio modules need a supported adapter; D-Bus BLE needs \texttt{python3-dbus}.
\end{itemize}

% ===========================================================================
\section{Learning resources}
For fundamentals, the IoT firmware-security canon: OWASP Firmware Security Testing
Methodology (FSTM), OWASP IoT Top 10, HackTricks (hardware/firmware), the Tarlogic
FSTM write-ups, and practitioner UART/SPI/\texttt{tar}/\texttt{nc} extraction
guides. For BLE: the LE Secure Connections vs legacy-pairing literature, and the
KNOB/BIAS/SweynTooth advisories. These informed the methodology implemented here.

% ===========================================================================
\section{Appendix --- command cheat-sheet}
\begin{table}[h]
\centering\footnotesize\intabletrue
\begin{tabularx}{\textwidth}{@{}L{6.5cm} X@{}}
\toprule
\textbf{Command} & \textbf{Purpose}\\
\midrule
\texttt{make api} / \texttt{make web} & Run backend / console (dev)\\
\texttt{pytest -q} & Run the backend test suite\\
\texttt{modprobe mac80211\_hwsim radios=3} & Stand up the no-RF Wi-Fi lab\\
\texttt{dd if=/dev/mtd0 | base64} & UART partition acquisition (per partition)\\
\texttt{tar cf - / | gzip | nc <host> 4444} & rootfs\_archive netcat transport\\
\texttt{flashrom -p linux\_spi:dev=/dev/spidev0.0} & SPI flash read-out\\
\texttt{C3I\_CVE\_ONLINE=false} & Use the offline CVE cache only\\
\bottomrule
\end{tabularx}
\end{table}

\vfill
\begin{center}\footnotesize\color{slate}
End of report · C3i Hub IoT Security Testbed · generated from LaTeX source
\end{center}

\end{document}
